\section{The \texttt{jml-ci} CLI}
\label{sec:cli}

The reference workflow is implemented as a Java CLI driver that the platform-specific workflow templates invoke. The CLI surfaces four subcommands:

\begin{lstlisting}[language={}]
jml-ci infer    --root <dir>  [--diff <base..head>]  [--cache-dir <dir>]
jml-ci verify   --root <dir>  [--diff <base..head>]  [--report <file>]
jml-ci embed    --jar <input.jar> --output <output.jar>  [--source <dir>]
jml-ci summary  --infer-report <file> --verify-report <file>
                [--format markdown|json] [--out <file>]
\end{lstlisting}

\paragraph{\texttt{infer}.}
For each {\tt .java} file in scope, look up its content hash in the cache; if hit, skip; if miss, run the inferrer and write the resulting JML annotations in place. Records the per-file outcome to a JSON report consumable by the \texttt{summary} subcommand. The {\tt --diff} flag restricts the scope to files changed in the supplied git diff; absent, the full root tree is processed.

\paragraph{\texttt{verify}.}
Discharge inferred specifications through OpenJML's Extended Static Checker. When {\tt OPENJML\_HOME} is unset, the subcommand fails open: it logs the absence and reports {\tt status: skipped} so the {\tt summary} subcommand can render the appropriate message. The {\tt --strict} flag converts a non-zero number of flagged clauses into a non-zero exit code; absent, exit is always zero.

\paragraph{\texttt{embed}.}
Run the inferrer over a source tree and embed the resulting specifications into a JAR via {\tt @JmlSpec} annotations (per Article~2). This is the deployment-time complement of {\tt infer}: whereas {\tt infer} writes JML into source files, {\tt embed} writes it into the compiled JAR so downstream consumers can read it without source access.

\paragraph{\texttt{summary}.}
Render a markdown summary of the most recent inference + verification run, suitable for posting as a PR comment. The {\tt --format json} variant emits machine-readable JSON for downstream tooling that wants the numeric outcomes.

\paragraph{Fail-open semantics throughout.}
Every subcommand exits with status zero by default, regardless of inference or verification outcomes. The opt-in {\tt --strict} flag is the only way to get a non-zero exit, and it is supplied by the platform-specific templates only when the consumer repo's CI configuration explicitly opts in. This is the design choice that makes the workflow safe to deploy in repos where formal-method tooling has previously been resisted on the grounds that it disrupts the merge train.

\paragraph{Test coverage.}
The CLI has a unit-test suite ({\tt JmlCiTest}, 7 tests) covering the help / no-args path, fail-open semantics on missing OpenJML, summary rendering with absent reports, content-sensitive cache key hashing, callee-spec sensitivity in the per-method cache key, and option-parsing for both flag and key-value forms. The full inferrer behaviour is exercised by the existing analyser-tier tests (101 tests) plus the embedder-tier tests (28 tests); the CLI's role is the workflow orchestration, not the inference, so the tests above are the right scope.
